\usepackage{tabularx,titlesec}
\usepackage{graphicx,color,fancyhdr,microtype}
\usepackage{booktabs,enumerate}
\usepackage{framed}
\usepackage{amsmath,amssymb,nth,mathtools,commath,nicefrac,siunitx,bigints,subcaption}
\usepackage{csquotes}

\usepackage[CJKbookmarks]{hyperref}
\hypersetup{
    colorlinks=false,
    %linkcolor=blue,
    %filecolor=magenta,      
    %urlcolor=cyan,
    %pdftitle={Overleaf Example},
    %pdfpagemode=FullScreen,
    }
    
\usepackage[a4paper,top=2.5cm,left=2cm,right=2cm]{geometry}
\usepackage{caption,subcaption} % use "subfigure"

\usepackage{xfrac} % use slanted fraction

\usepackage{physics} % to use bra, ket

\usepackage[utf8]{inputenc} % tell latex it's under utf8

\pagestyle{fancy}
\lfoot{\includegraphics[scale=.05]{wise_logo_2016.pdf}}
\rfoot{\includegraphics[scale=.4]{fine_cornell_logo_red_cmyk.eps}}

% bibliography
\usepackage[
	backend=bibtex,
	bibstyle=phys,
	citestyle=numeric-comp,
	articletitle=true,
	biblabel=brackets,
	sorting=none
	]{biblatex}
\addbibresource{reference.bib}

\newcommand*{\Exp}{\mathrm{e}}
\newcommand{\sgn}{\operatorname{sgn}}
\DeclarePairedDelimiter\floor{\lfloor}{\rfloor}
\newcommand{\pvec}[1]{\vec{#1}\mkern2mu\vphantom{#1}}

% Line break if equations are too long
\allowdisplaybreaks[1]

% Differential "d" in the integral. It needs amsmath
\newcommand*\diff{\mathop{}\!\mathrm{d}}
\newcommand*\Diff[1]{\mathop{}\!\mathrm{d^#1}}

\usepackage[version=3]{mhchem}

% MATLAB code
\usepackage{listings,lstautogobble,verbatim}
\definecolor{mygreen}{RGB}{28,172,0} % color values Red, Green, Blue
\definecolor{mylilas}{RGB}{170,55,241}
\definecolor{codebg}{RGB}{237,245,207}
\lstset{language=matlab,
	backgroundcolor=\color{codebg},
    basicstyle=\footnotesize,
    breaklines=true,%
    morekeywords={matlab2tikz},
    keywordstyle=\color{blue},
    morekeywords=[2]{1}, keywordstyle=[2]{\color{black}},
    identifierstyle=\color{black},
    stringstyle=\color{mylilas},
    commentstyle=\color{mygreen},
    showstringspaces=false,%without this there will be a symbol in the places where there is a space
    numbers=left,%
    numberstyle={\tiny \color{black}},% size of the numbers
    numbersep=9pt, % this defines how far the numbers are from the text
    emph=[1]{for,end,break},emphstyle=[1]\color{red}, %some words to emphasise
    autogobble=true,
}

\mathchardef\mhyphen="2D % use hyphen in the math mode; "2D is the ASDII